\documentclass[11pt,a4paper]{article}
\usepackage[utf8]{inputenc}
\usepackage[T1]{fontenc}
\usepackage[spanish,es-nodecimaldot]{babel}
\usepackage{lmodern}
\usepackage[margin=2.4cm,headheight=14pt]{geometry}
\usepackage{booktabs,tabularx,array}
\usepackage{enumitem}
\usepackage{xcolor}
\usepackage{fancyhdr}
\usepackage{hyperref}

\definecolor{verde}{HTML}{174A45}
\definecolor{gris}{HTML}{475569}
\hypersetup{colorlinks=true,linkcolor=verde,urlcolor=verde}
\setlength{\parskip}{5pt}
\setlength{\parindent}{0pt}
\setlist{leftmargin=*,itemsep=2pt,topsep=3pt}
\pagestyle{fancy}
\fancyhf{}
\lhead{Programación y bases de datos}
\rhead{Avance 2 -- Proyecto integrador}
\cfoot{\thepage}
\renewcommand{\headrulewidth}{0.4pt}
\newcommand{\completar}[1]{\textit{\textcolor{gris}{#1}}}

\begin{document}
\hypersetup{pageanchor=false}

\begin{titlepage}
\centering
\vspace*{2.0cm}
{\Large\bfseries UNIVERSIDAD NACIONAL DE LOJA\par}
\vspace{0.3cm}
{\large Facultad Agropecuaria y de Recursos Naturales Renovables\par}
{\large Carrera de Ingeniería Ambiental\par}
\vspace{1.3cm}
{\color{verde}\rule{0.82\linewidth}{1.3pt}\par}
\vspace{0.55cm}
{\huge\bfseries \completar{Título del proyecto: tema y zona}\par}
\vspace{0.45cm}
{\Large Avance 2\par}
\vspace{0.55cm}
{\color{verde}\rule{0.82\linewidth}{1.3pt}\par}
\vfill
\begin{tabular}{rl}
\textbf{Asignatura:} & Programación y bases de datos\\
\textbf{Grupo:} & \completar{Número de grupo}\\
\textbf{Integrantes:} & \completar{Nombres y apellidos completos}\\
 & \completar{Añada filas para todos los integrantes}\\
\textbf{Período lectivo:} & octubre 2026 -- febrero 2027\\
\end{tabular}
\vspace{1cm}
\end{titlepage}
\hypersetup{pageanchor=true}

\section{Problema de investigación}
\subsection{Descripción del problema}
Describan la situación concreta relacionada con el tema y la zona asignada. Organicen el texto así:
\begin{enumerate}
  \item Situación territorial o ambiental que se analizará.
  \item Vacío de información o dificultad específica.
  \item Relevancia del problema, sin atribuir a la base efectos que no mide.
  \item Oportunidad de responderlo con el archivo, variables y análisis del grupo.
\end{enumerate}

\completar{Redacte aquí el problema con datos de contexto y citas cuando corresponda. Sustituya estas instrucciones antes de entregar.}

\subsection{Pregunta de investigación}
Formulen una pregunta única, clara y respondible con cálculos sobre la base. Incluyan variable o fenómeno, zona y período 2019--2024.

\textbf{Pregunta:} \completar{Escriba aquí la pregunta de investigación.}

\subsection{Hipótesis}
Escriban una afirmación contrastable que adelante una respuesta. Indiquen variable, zona, período, indicador y criterio cuantitativo. No usen expresiones sin definir como «mucho», «normal» o «afecta».

\textbf{Hipótesis:} \completar{Escriba aquí una afirmación medible y verificable con la base.}

\section{Objetivos}
\subsection{Objetivo general}
Debe responder directamente la pregunta e indicar qué se analizará, dónde, con qué datos y en qué período.

\textbf{Objetivo general:} \completar{Escriba un objetivo general con verbo en infinitivo.}

\subsection{Objetivos específicos}
Incluyan máximo tres objetivos. Cada uno debe ser medible, aportar al objetivo general y producir una salida verificable; redacten objetivos, no tareas.
\begin{enumerate}
  \item \completar{Objetivo específico 1 y su salida verificable.}
  \item \completar{Objetivo específico 2 y su salida verificable.}
  \item \completar{Objetivo específico 3 (opcional) y su salida verificable.}
\end{enumerate}

\section{Metodología}
\subsection{Zona de estudio y fuente de datos}
Indiquen cantón, zona, período, nombre exacto del CSV y unidad de observación. Expliquen cómo comprobarán que la base corresponde a su asignación. Las bases ya vienen separadas por grupo: no las unan ni vuelvan a dividir. La muestra didáctica describe únicamente las celdas entregadas, no una estimación oficial de todo el cantón. El archivo CSV es de solo lectura.

\textbf{Grupo y tema:} \completar{Número de grupo y bloque temático.}\\
\textbf{Zona:} \completar{Norte, sur o Loja completo; explique la regla aplicable.}\\
\textbf{Archivo asignado:} \completar{Nombre exacto del CSV.}\\
\textbf{Período y observación:} \completar{2019--2024; celda de 500 m en un mes, u otra precisión pertinente.}

\subsection{Variables seleccionadas}
Completen con los nombres exactos de columnas del CSV asignado y consulten significado y unidad en \texttt{DICCIONARIO\_VARIABLES.csv}. Incluyan solo variables que se usarán para responder la pregunta y distingan variables analíticas de controles de calidad.

\small
\begin{tabularx}{\linewidth}{>{\raggedright\arraybackslash}p{4.0cm} X >{\raggedright\arraybackslash}p{2.1cm} >{\raggedright\arraybackslash}p{1.7cm} >{\raggedright\arraybackslash}p{2.2cm}}
\toprule
\textbf{Variable (columna)} & \textbf{Significado} & \textbf{Unidad de medida} & \textbf{Tipo} & \textbf{Uso en el análisis}\\
\midrule
\completar{Nombre exacto} & \completar{Definición según el diccionario} & \completar{Unidad} & \completar{Tipo} & \completar{Analítica o control}\\
\completar{Nombre exacto} & \completar{Definición según el diccionario} & \completar{Unidad} & \completar{Tipo} & \completar{Analítica o control}\\
\completar{Nombre exacto} & \completar{Definición según el diccionario} & \completar{Unidad} & \completar{Tipo} & \completar{Analítica o control}\\
\bottomrule
\end{tabularx}
\normalsize

\subsection{Exploración preliminar y procedimiento}
Registren qué se hizo y qué se encontró al revisar el CSV, sin modificarlo:
\begin{itemize}
  \item Filas y columnas revisadas: \completar{indique conteos observados}.
  \item Filas de ejemplo consultadas: \completar{indique cómo se seleccionaron; adjunte una tabla breve si es necesario}.
  \item Tres hechos observados: \completar{describa hechos verificables en variables concretas}.
  \item Control de periodo, zona, valores faltantes y disponibilidad: \completar{explique qué verificó}.
\end{itemize}

Expliquen, en correspondencia con cada objetivo, las columnas de entrada, filtros o criterios de calidad, cálculo previsto y salida verificable. Diferencien las acciones realizadas de las que quedan planificadas. No incorporen datos de 2025 ni sumen ventanas móviles solapadas como si fueran independientes.

\subsection{Fuentes y referencias}
Citen las fuentes usadas para contextualizar el problema y documenten la base y el diccionario compartidos. Mantengan correspondencia entre citas en el texto y referencias; apliquen APA 7.

\begin{itemize}
  \item \completar{Autor o institución. (Año). Título. URL o datos de publicación.}
  \item \completar{Agregue las fuentes realmente consultadas.}
\end{itemize}

\end{document}